\chapter{Colinéarité, déterminant et bases}
\section*{Exercice 9.1 -- Déterminants}
Le premier vaut $2\times15-5\times6=0$ : colinéarité, facteur $3$. Le second vaut $14-15=-1\ne0$ : les vecteurs ne sont pas colinéaires.

\section*{Exercice 9.2 -- Alignement}
$\overrightarrow{AB}=(2,3)$ et $\overrightarrow{AC}=(4,6)=2\overrightarrow{AB}$. Comme $A\ne B$, cela prouve que $C$ appartient à $(AB)$.

\section*{Exercice 9.3 -- Paramètre}
Le déterminant vaut $3t-12$. Il est nul exactement pour $t=4$. Alors $(4,6)=2(2,3)$.

\section*{Exercice 9.4 -- Coordonnée nulle}
Le déterminant vaut $0\times(-7)-4\times0=0$. On peut aussi écrire $(0,-7)=(-7/4)(0,4)$ ; on divise par la coordonnée non nulle, pas par $0$.

\section*{Exercice 9.5 -- Dans une autre base}
Le système est $x+2y=7$, $2x-y=1$. De $y=2x-1$, on tire $5x=9$, donc $x=9/5$ et $y=13/5$. La base est valide car son déterminant est $-1-4=-5\ne0$.

\section*{Exercice 9.6 -- Parallèles ou confondues}
Les directeurs $(2,1)$ et $(4,2)$ sont colinéaires. Mais $C$ n'appartient pas à $(AB)$ : $\det((2,1),(0,3))=6\ne0$. Les droites sont parallèles distinctes.
